\subsection{Variantenkonzept}\label{subsec:varianten}

Aus den in \autoref{sec:personas} beschriebenen Personas und den daraus abgeleiteten Nutzungskontexten wurden zwei konzeptionell unterschiedliche Varianten abgeleitet. Sie unterscheiden sich nicht nur in der Bildschirmgröße, sondern auch in Navigationsmodell, Informationsarchitektur und Interaktionslogik.

\begin{table}[H]
    \centering
    \scriptsize
    \begin{tabular}{@{}p{2.7cm}p{7.0cm}p{6.0cm}@{}}
        \toprule
        \textbf{Attribut} & \textbf{Variante 1 -- Desktop-Variante} & \textbf{Variante 2 -- Mobil-Variante} \\
        \midrule
        Zielgruppe & Registrierte Stammnutzer:innen, z.~B. Dr.~med.~Weber als Mehrheitsnutzungskontext & Anonyme Empfänger:innen, z.~B. Markus Fischer; zusätzlich registrierte mobile Nutzer:innen im Multi-Device-Setting \\
        \addlinespace
        Geräte & Desktop-PC, iPad im Querformat & Tablet im Hochformat, Smartphones \\
        \addlinespace
        Primärnavigation & Dauerhaft sichtbare \emph{Drawer}-Navigation & \textit{Off canvas}-Navigation mit Hamburger-Menü \parencite[Kap.~15.1.10]{erlhoferWebsiteKonzeptionUndRelaunch2017} \\
        \addlinespace
        Layout & \emph{Vier-Spalten-Layout} mit Drawer, SideNav, Master und Detail & \emph{Ein-spaltiges Layout} mit Overlay-Strategie für Detail-Informationen \\
        \addlinespace
        Sekundärnavigation & Vertikale Sekundärnavigation, dauerhaft sichtbar & Angepasste horizontale Filterleiste, kontextuell eingebunden \\
        \addlinespace
        Touch-Targets & Standardgröße für Maus- und Tastaturbedienung & Optimiert für Finger-Bedienung \\
        \bottomrule
    \end{tabular}
    \caption[Varianten-Vergleich Desktop vs.\ Mobil]{Konzeptionelle Unterschiede der Desktop- und Mobil-Variante. Quelle: Eigene Darstellung.}
    \label{tab:varianten_vergleich}
\end{table}

Beide Varianten basieren auf demselben UI-Framework und derselben Backend-Schnittstelle. Die Umschaltung erfolgt über Breakpoints, sodass Geräte wie Tablets je nach Viewport und Ausrichtung die jeweils passende Variante nutzen können. Die Breakpoint-Steuerung dient dabei nicht nur der technischen Skalierung, sondern der Übergabe zwischen zwei eigenständig konzipierten Nutzungserlebnissen.

Die konzeptionellen Unterschiede -- insbesondere Navigationsmuster, Layout-Strategien und Touch-Targets -- werden in \autoref{sec:skizzen_wireframes} visuell konkretisiert. Die finalen Mockups beider Varianten sind in \autoref{sec:mockups} dokumentiert.

\paragraph{Prototyping- und Fidelity-Strategie}
Zur Ausarbeitung beider Varianten wurde die Designlösung schrittweise von niedriger zu hoher Fidelität konkretisiert. Ziel war es, zunächst Struktur- und Navigationsentscheidungen unabhängig von visuellen Details zu prüfen und anschließend in eine realitätsnahe, klickbare Umsetzung zu überführen. \autoref{fig:fidelitaetshierarchie} zeigt die eingesetzten Stufen von manuellen Skizzen bis zur lauffähigen Webanwendung.

\begin{figure}[ht!]
    \centering
    \begin{tikzpicture}[
        node distance=0.8cm,
        box/.style={
            rectangle,
            draw,
            rounded corners,
            fill=blue!10,
            minimum width=3.5cm,
            minimum height=1.5cm,
            align=center,
            font=\small
        },
        arrow/.style={->, >=latex, thick}
    ]
        \node[box] (a) {Manuelle Skizzen\\\small iPad\\(Low-Fidelity)};
        \node[box] (b) [right=of a] {Digitale Wireframes\\\small Figma\\(Mid-Fidelity)};
        \node[box] (c) [right=of b] {Implementierung\\\small Vaadin Flow\\(High-Fidelity)};
        \node[box] (d) [right=of c] {Lauffähige\\\small Webanwendung};

        \draw[arrow] (a) -- (b);
        \draw[arrow] (b) -- (c);
        \draw[arrow] (c) -- (d);
    \end{tikzpicture}
    \caption[Fidelity-Hierarchie]{Fidelity-Hierarchie im Redesign-Prozess. Quelle: Eigene Darstellung.}
    \label{fig:fidelitaetshierarchie}
\end{figure}

Die Low- und Mid-Fidelity-Stufen werden in \autoref{sec:skizzen_wireframes} dokumentiert. Die High-Fidelity-Umsetzung erfolgte direkt in Vaadin Flow und bildet die Grundlage der in \autoref{sec:mockups} gezeigten Mockups sowie des in \autoref{sec:usability_test_evaluation} evaluierten Prototyps.
